\documentclass[a4paper,11pt]{article}

% Front matter
\usepackage[T1]{fontenc}
\usepackage[latin1]{inputenc}
\usepackage[francais,english]{babel}
\selectlanguage{english}
\usepackage{amsmath,amsthm,amssymb}
\usepackage{fancyhdr}
\usepackage[top=2.5cm, bottom=2.25cm, left=2.5cm, right=2.5cm]{geometry}
\usepackage{makecell}
\usepackage{graphicx}
\usepackage{tabularx}
\usepackage{makeidx}
\usepackage{minitoc}
\usepackage{hyperref}
\usepackage{symbmath}
\usepackage[ruled]{algorithm2e}
\usepackage{caption}
\usepackage{subcaption}

\renewcommand{\algorithmcfname}{Algorithm}

% ----------------
% New environments
% ----------------

\newtheorem{Defi}{Definition}%[section]
\newtheorem{Not}{Notation}
\newtheorem{Lem}{Lemma}%[section]
\newtheorem{Hypo}{Hypothesis}
\newtheorem{Theo}{Theorem}%[section]
\newtheorem{Prop}{Proposition}%[section]
\newtheorem{Coro}{Corollary}%[section]
\newtheorem{Rem}{Remark}
\newtheorem{Pb}{Problem}

% -----
% Title
% -----

\title{Biharmonic problem solved by a discontinuous Galerkin finite element method}

% ---------------------
% Begin of the document
% ---------------------

\begin{document}

\maketitle

\section{Problem statement}

The biharmonic equation appears in various physics problems as for instance in acoustic, fluid mechanics among others. The aim of this numerical exercise is to solve the biharmonic equation with a mixed formulation using a discontinuous Galerkin finite element method.

We consider the following problem

\begin{Pb}
  Find $u$, defined in $\Omega$, such that
  \begin{eqnarray}
    -\Delta^2u=f\mbox{ in }\Omega,\\
     u=g_1\mbox{ on }\partial\Omega,\\
     \frac{\partial u}{\partial n}=g_2\mbox{ on }\partial\Omega,
  \end{eqnarray}
\end{Pb}

In this problem, $f$, $g_1$ and $g_2$ are three given functions. We assume that $f$, $g_1$ and $g_2$ are sufficiently regular so that a unique solution $u\in H^4(\Omega)$ exists.

In order to use a mixed formulation, one transforms the previous problem as follows
\begin{Pb}
  Find $u$ and $v$, defined in $\Omega$, such that
  \begin{eqnarray}
    -\Delta u+v=0\mbox{ in }\Omega,\\
    -\Delta v=f\mbox{ in }\Omega,\\
     u=g_1\mbox{ on }\partial\Omega,\\
     \frac{\partial u}{\partial n}=g_2\mbox{ on }\partial\Omega,
  \end{eqnarray}
\end{Pb}

\section{Discontinuous finite element formulation}

Let $\mathcal{T}_h$ be a mesh of the domain $\Omega$ with meshsize $h$ formed by finite elements $K$ such as $\mathcal{T}_h=\{K\}$ and $\overline{\Omega}=\bigcup_{K\in\mathcal{T}_h}\overline{K}$. Furthermore, we consider a subset of faces $F$ formed by interfaces between two distint mesh elements $K_1$ and $K_2$ such as $F=\partial T_1\cap\partial T_2$ and boundary faces given by $F=\partial T\cap\partial \Omega$. If $\mathcal{F}^i_h$ gathers all interfaces and $\mathcal{F}^b_h$ all boundary faces, the set of faces is $\mathcal{F}_h=\mathcal{F}^i_h\cup\mathcal{F}^b_h$.

The discontinuous finite element space is defined by
\begin{equation}
  X_h=\{v_h\in L^2(\Omega);\ v_{h|K}\in P_k,\ \forall\ K\in \mathcal{T}_h\},
\end{equation}
with $k$ the polynomial degree. This space belongs to the broken Sobolev space
\begin{equation}
  H^1(\mathcal{T}_h)=\{v\in L^2(\Omega);\ v_{|K}\in H^1(K),\ \forall\ K\in \mathcal{T}_h\}.
\end{equation}

The discretization is established according to developments used for harmonic equation \cite[Chap. 4]{DIPIETROERN2012}. The following bilinear forms are introduced
\begin{eqnarray}
  a_h(u,w)=\int_{\Omega}\nabla_hu\cdot\nabla_hwdV-\sum_{F\in\mathcal{F}_h}\int_{F}\average{\frac{\partial u}{\partial n}}
           \jump{w}dS-
           \sum_{F\in\mathcal{F}^i_h}\int_{F}\average{\frac{\partial w}{\partial n}}\jump{u}dS,\\
  J_h(u,q)=\sum_{F\in\mathcal{F}_h}\eta_{|F}\int_{F}\jump{u}\jump{q}dS,
\end{eqnarray}
in which $\eta_{|F}$ is given by \cite{SHAHBAZI2005}
\begin{equation}
  \eta_{|F}=\frac{(k+1)(k+d)\varpi_{|F}}{d},
\end{equation}
with $\varpi_{|F}$
\begin{equation}
  \varpi_F=\left\lbrace
      \begin{array}{lr}
        \frac{\mathrm{meas}(\partial K)}{\mathrm{meas}(K)}&\mbox{ for }F\in K\cap\partial\Omega,\\
        \max\left(\frac{\mathrm{meas}(\partial K_+}{\mathrm{meas}(K_+)},\frac{\mathrm{meas}(\partial K_-}{\mathrm{meas}(K_-)}\right)&
        \mbox{ for }F\in K_-\cap K_+.
      \end{array}
      \right.
\end{equation}
Notice that $\varpi_F$ scales as $h_{|F}^{-1}$ with $h_{|F}$ the local length scale of the current face $F$.

We recall also that the average and jump at each face between two elements $K_1$ and $K_2$ are given by
\begin{eqnarray}
  \average{u}=\frac{1}{2}(u|_{K_1}+u|_{K_2}),\\
  \jump{u}=u|_{K_1}-u|_{K_2}.
\end{eqnarray}

Linear forms are also defined as follows
\begin{eqnarray}
  l_{h1}(w)=\sum_{F\in\mathcal{F}^b_h}\int_{F}\left(g_2w-g_1\frac{\partial w}{\partial n}\right)dS,\\
  l_{h2}(q)=\int_{\Omega}fqdV+\sum_{F\in\mathcal{F}^b_h}\int_{F}\alpha_{|F}g_1qdS.
\end{eqnarray}

So, the discontinuous problem is formulated under the coupled form as follows
\begin{Pb}
  Find $u_h$, $v_h\in X_h$ such that
  \begin{eqnarray}
    a_h(w_h,u_h)+(v_h,w_h)=l_{h1}(w_h),\mbox{ for }w_h\in X_h,\\
    J_h(u_h,q_h)+a_h(v_h,q_h)=l_{h2}(q_h),\mbox{ for }q_h\in X_h.
  \end{eqnarray}
\end{Pb}

\section{Numerical example}

The previous problem has been implemented in C++ Rheolef library \cite{RHEOLEF2015,RHEOLEFDG2015}. The numerical example is chosen in order to have an exact solution useful to determine the numerical errors. We consider a domain $\Omega=(0,1)\times(0,1)$. The exact solution of $u$ is given by
\begin{equation}
  u(x,y)=\beta x^2(1-x)^2y^2(1-y)^2.
\end{equation}
Consequently, the boundary conditions becomes homogeneous, i.e. $g_1=g_2=0$, and the source term $f$ is given by
\begin{equation}
  f(x,y)=-8\beta\left\lbrace3\left[x^2(1-x)^2+y^2(1-y)^2\right]+(1-6x+6x^2)(1-6y+6y^2)\rbrace\right.
\end{equation}

The numerical example has been achieved in regular square with mesh size varying in the range $h\in[1/16;1/64]$. The coefficient $\beta$ is set at $100$. Figures \ref{subfig:uvsxy05} and \ref{subfig:vvsxy05} depict the solution found for $h=1/16$ and $1/32$ and for two polynomial degrees equal to $2$ and $3$. Agreement with the exact solution is very well established.

\begin{figure}
  \centering
  \begin{subfigure}{0.47\textwidth}
    \centering
    \caption{$u$}
    \includegraphics[width=0.98\textwidth]{uvsxy05.pdf}
    \label{subfig:uvsxy05}
  \end{subfigure}
  \begin{subfigure}{0.47\textwidth}
    \centering
    \caption{$v$}
    \includegraphics[width=0.98\textwidth]{vvsxy05.pdf}
    \label{subfig:vvsxy05}
  \end{subfigure}
  \caption{$u$ and $v$ vs $x$ along the horizontal axis $y=1/2$ for the two mesh sizes and two polynomial degrees.}
  \label{fig:uvvsxy05}
\end{figure}

In order to be more quantitative on the scheme accuracy, we determine the errors in two norms. The first is the $L^2$ norm defined as usual by
\begin{equation}
  ||u||_{L^2}^2=\int_{\Omega}u^2dV.
\end{equation}
The second norm is a mesh-dependent norm defined as follows \cite{GUDIETAL2008}
\begin{equation}
  |||u|||^2=\int_{\Omega}\nablagras_hu\cdot\nablagras_hudV+\sum_{F\in\mathcal{F}_h}\int_{F}\frac{|h_{|F}|}{k^2}
  \average{\frac{\partial u}{\partial n}}^2 dS+
  \sum_{F\in\mathcal{F}_h}\int_{F}\frac{k^2}{|h_{|F}|}\jump{u}^2dS.
\end{equation}

Figures \ref{subfig:erroruP2dP3dbihar} and \ref{subfig:errorvP2dP3dbihar} plot the errors following the two previous norms obtained for two polynomial degrees equal to $P2d$ and $P3d$. The slope of convergence has been determined using a regression routine. The exponent of the error $n$ is given for each curve in the previous figures. These error estimations show that the scheme is optimal-order in error bounds.

\begin{figure}
  \centering
  \begin{subfigure}{0.47\textwidth}
    \centering
    \caption{$u$}
    \includegraphics[width=0.98\textwidth]{erroruP2dP3dbihar.pdf}
    \label{subfig:erroruP2dP3dbihar}
  \end{subfigure}
  \begin{subfigure}{0.47\textwidth}
    \centering
    \caption{$v$}
    \includegraphics[width=0.98\textwidth]{errorvP2dP3dbihar.pdf}
    \label{subfig:errorvP2dP3dbihar}
  \end{subfigure}
  \caption{Errors on $u$ and $v$ vs $h$ for two polynomial degrees.}
  \label{fig:uvvsxy05}
\end{figure}

% ------------
% Bibliography
% ------------

\bibliographystyle{plain}
\bibliography{bibliojabref}

\end{document}
